\section*{Chapitre \ref{ch:nw:capturing-and-monitoring-the-network} --- Capture et supervision du réseau}

\begin{solution}{exo:nw:capturing-and-monitoring-the-network:1}
Non~: le champ des nanosecondes doit être inférieur à $10^9$. Un lecteur indulgent reporterait la seconde et lirait 34\,201~s et 500\,000\,321~ns,
soit 09:30:01.500\,000\,321~; un lecteur strict devrait rejeter l'enregistrement, parce que l'écrivain est défaillant.
\end{solution}

\begin{solution}{exo:nw:capturing-and-monitoring-the-network:2}
$16\,771 - 12\,345 = \qty{4\,426}{\nano\second}$, à partir de la première copie (ligne~A). Lorsque la copie de la ligne~A est perdue, la stratégie voit le paquet sur la ligne~B
\qty{3}{\micro\second} plus tard, et son ordre serait parti environ \qty{3}{\micro\second} plus tard~: mesurée à partir de la copie B, la latence wire-to-wire
est la même, et les \qty{3}{\micro\second} appartiennent à la ligne perdue.
\end{solution}

\begin{solution}{exo:nw:capturing-and-monitoring-the-network:3}
Chaque sens transporte $4 \times 10^9 / (220 \times 8) \approx 2.27$ million de trames par seconde (200 octets, plus 20 de préambule et d'espace inter-trames sur
le fil)~; chacune est stockée sur $200 + 16 + 16 = 232$ octets~: environ 527~Mo par seconde et par sens, soit 1.05~Go par seconde pour la liaison.
\end{solution}

\begin{solution}{exo:nw:capturing-and-monitoring-the-network:4}
La règle affecte un ordre au dernier déclencheur qui le précède. Elle est fausse dès qu'un autre déclencheur arrive entre le véritable déclencheur et
l'ordre, ce qui se produit avec une probabilité qui croît avec le taux de déclenchement multiplié par la latence de la société~: des déclencheurs
fréquents, ou une société lente, placent davantage de déclencheurs dans la fenêtre.
\end{solution}

\begin{solution}{exo:nw:capturing-and-monitoring-the-network:5}
$99\,055 \times 0.002^2 \approx 0.4$~: n'en voir aucun est ce que prédit l'indépendance. Un paquet perdu sur les deux lignes a une cause commune.
\end{solution}

\begin{solution}{exo:nw:capturing-and-monitoring-the-network:6}
La médiane~: les deux passages de la carte (1.6~\unit{\micro\second} au total, une meilleure carte ou le contournement du noyau sur un hôte réglé) et l'anneau
(un consommateur en attente active sur un cœur isolé). Le 99\up{e}~centile~: l'anneau, dont la queue provient de l'ordonnancement~; isoler ses cœurs coûte
des cœurs et du réglage (livre~13, chapitre~13), et non de l'argent.
\end{solution}

\begin{solution}{exo:nw:capturing-and-monitoring-the-network:7}
L'erreur de la règle naïve tombe de 47~\% à 25~\%~: avec une société quatre fois plus rapide, moins de déclencheurs arrivent entre un déclencheur et son ordre. Elle ne disparaît pas,
car les \omterm{def:nw:capturing-and-monitoring-the-network:trigger}{paquets déclencheurs} arrivent encore par rafales plus rapprochées que la latence de la société.
\end{solution}

\begin{solution}{exo:nw:capturing-and-monitoring-the-network:8}
Les horodatages de la carte ne voient rien de ce qui se trouve hors de la carte~: la fibre de la cage et les équipements des deux côtés, que le tap observe. Ils reposent aussi sur
l'horloge de la carte et sur la façon dont elle définit ses instants de réception et d'émission. C'est une bonne mesure interne à l'hôte~; ce n'est pas la latence wire-to-wire
de la société aux points de remise.
\end{solution}

\begin{solution}{pb:nw:capturing-and-monitoring-the-network:1}
\begin{enumerate}
\item $437 + 756 + 158 + 97 + 47 + 33 = \qty{1\,528}{\nano\second}$.
\item $1\,528 + 263 + 295 + 798 + 801 = \qty{3\,685}{\nano\second}$.
\item Les étapes sont dissymétriques vers la droite~: la médiane d'une somme de variables dissymétriques se situe au-dessus de la somme de leurs médianes (ici
  \qty{4\,386}{\nano\second} dans le budget).
\item $4.4 - 1.5 = \qty{2.9}{\micro\second}$~: soit les deux tiers de la médiane.
\item $31.9 / 38.5 = 83~\%$.
\item Un centile élevé de la somme est dominé par l'étape dont la queue est la plus lourde~; les queues des autres étapes coïncident rarement avec elle.
\item Un 99.9\up{e}~centile de \qty{3.5}{\milli\second}~: le consommateur a été désordonnancé sur des cœurs partagés.
\item Des cœurs isolés et un consommateur en attente active, qui retirent l'ordonnanceur de l'anneau.
\item Le tap voit la trame à l'instant où la société la reçoit de la plateforme~; l'horodatage de la carte intervient après le réseau de la cage, il manque donc
  une partie du chemin, et l'horloge de la carte est une source de plus à laquelle il faut se fier.
\item Sous-estimé~: la règle naïve associe les ordres lents à des paquets plus tardifs, et donne un 99\up{e}~centile de \qty{27}{\micro\second} au lieu de 38.
\item L'identifiant d'ordre client, ou un champ que la plateforme renvoie inchangé.
\item Des centiles wire-to-wire calculés à partir de la capture chaque minute, comparés à une référence.
\item \textbf{Résultat nommé.} Médiane wire-to-wire \qty{4.4}{\micro\second} et 99\up{e}~centile \qty{38}{\micro\second}~; \qty{2.9}{\micro\second} de la
  médiane se situent hors du programme~; l'anneau entre le gestionnaire et la stratégie porte la queue.
\item D'environ \qty{0.9}{\micro\second} (la médiane du budget passe de 4.39 à 3.47).
\item Le décodage (p99 \qty{2.8}{\micro\second}), suivi de près par les deux passages de la carte, à \qty{2}{\micro\second} chacun.
\item Des horodatages matériels au port de la carte des deux côtés, comparés à ceux des taps, ou une capture de chaque côté du serveur.
\item C'est ce que la société a effectivement envoyé et reçu, et à quel moment, indépendamment des journaux de tout programme~: l'enregistrement qu'une plateforme, un client ou un
  régulateur réclame.
\item Environ 38~Mo par seconde de capture, soit 138~Go par heure.
\item Un mois de synthèses (centiles de latence, trous, rafales, alarmes) et des trames du transfert des ordres~; une journée de chaque trame de données de marché,
  sauf si un incident en demande davantage.
\item Tout ce qui se trouve hors du programme, daté par une horloge que le programme ne peut pas modifier.
\end{enumerate}
\end{solution}

\begin{solution}{iq:nw:capturing-and-monitoring-the-network:1}
Placer des taps sur le transfert des données de marché et sur celui des ordres, horodater les deux sur le fil avec une même horloge disciplinée, associer chaque ordre à son
\omterm{def:nw:capturing-and-monitoring-the-network:trigger}{paquet déclencheur}, puis soustraire. Les journaux d'un programme ne voient ni le réseau, ni la carte, ni le noyau des deux côtés, et utilisent l'horloge de l'hôte.
\iqlookfor{taps, une seule horloge et l'appariement.}
\end{solution}

\begin{solution}{iq:nw:capturing-and-monitoring-the-network:2}
Sur le fil, au premier octet de la trame, par le dispositif de prise de copie, transporté avec la copie (un suffixe). Tout ce qui intervient plus tard ajoute
la file d'attente du réseau d'agrégation et l'horloge du serveur de stockage à chaque mesure.
\iqlookfor{horodatage au tap, et non au stockage.}
\end{solution}

\begin{solution}{iq:nw:capturing-and-monitoring-the-network:3}
Faire en sorte que la stratégie porte le numéro de séquence du déclencheur (ou un condensat) dans l'ordre, dans un champ que la capture peut lire. Les heuristiques, comme le
dernier déclencheur précédant l'ordre, échouent lors des rafales, c'est-à-dire précisément quand la latence compte.
\iqlookfor{instrumenter l'ordre plutôt que de deviner.}
\end{solution}

\begin{solution}{iq:nw:capturing-and-monitoring-the-network:4}
Les rejets en sortie sur des ports qui ne devraient jamais être saturés, les trous par ligne et sur les deux lignes, les reprises par instantané, le verrouillage de l'horloge et ses décalages, et
les centiles wire-to-wire comparés à une référence~: chacun désigne un point du chemin.
\iqlookfor{alarmes liées à des mécanismes, et non un tableau de bord de tout.}
\end{solution}

\begin{solution}{iq:nw:capturing-and-monitoring-the-network:5}
Le décomposer~: comparer les horodatages par étape de la capture (taps, horodatages de carte s'ils sont conservés) avec le budget~; vérifier si le marché est devenu plus actif
(rafales, débits de messages), si une horloge a dérivé (maintien, nouvelle asymétrie), si un hôte a changé (noyau, micrologiciel, nouveau processus sur un
cœur partagé), et si la plateforme a modifié son flux. Examiner à quelle seconde cela a commencé.
\iqlookfor{décomposition par étape et liste des causes hors déploiement.}
\end{solution}

\begin{solution}{iq:nw:capturing-and-monitoring-the-network:6}
Analyse des trous et des pertes, détection des \omterm{def:nw:networking-for-trading:burst}{micro-rafales}, reconstruction du carnet pour vérifier le gestionnaire, reconstitution de ce que la société a envoyé pendant un incident,
vérification des délais d'accusé de réception de la plateforme, et preuves pour les litiges et les régulateurs.
\iqlookfor{la capture comme preuve et comme seconde implémentation indépendante.}
\end{solution}
